\capitulo{1}{Introducción}

\section{Contexto del proyecto}

Cuando se diseña una base de datos, antes de decidir su implementación en un modelo lógico concreto, es necesario representar de forma conceptual la información del dominio y las relaciones existentes entre sus elementos. El modelo Entidad-Relación, propuesto por Chen~\cite{chen1976}, permite realizar esta representación mediante entidades, atributos e interrelaciones, facilitando el análisis del problema antes de abordar su transformación posterior.

Este Trabajo Fin de Grado se sitúa en ese contexto, pero no aborda el desarrollo de una herramienta desde cero. El proyecto se centra en la evolución de UBU E-R App, anteriormente denominada Draw E-R App, una aplicación web para la creación de diagramas Entidad-Relación desarrollada con React y mxGraph a partir del trabajo previo de Rubén Maté Iturriaga~\cite{mate2024drawERApp} y del repositorio original de la aplicación~\cite{mate2024drawERRepo}. La aplicación heredada ya ofrecía una base funcional, pero requería análisis, estabilización y ampliación para mejorar la coherencia entre el diagrama editado en el lienzo, el modelo interno, las validaciones, la transformación relacional y la generación de SQL.

A fin de ofrecer una primera visión de la aplicación, la Figura~\ref{fig:editor_general} muestra una vista general del editor, donde se aprecia el lienzo de trabajo, el panel de acciones y la representación visual de un diagrama Entidad-Relación sencillo.

\imagen{editor_general}{Vista general del editor web con un diagrama Entidad-Relación sencillo.}{0.85}

\section{Motivación}

La motivación de este trabajo surge de una necesidad práctica detectada en el aprendizaje del diseño de bases de datos. Aunque existen numerosas herramientas para crear diagramas, muchas están orientadas al modelado con UML, a diagramas físicos de bases de datos o a notaciones de tipo \textit{crow's foot}. Algunas incluso se presentan como herramientas para diagramas E/R, aunque no sigan de forma fiel la notación utilizada habitualmente en la docencia del modelo Entidad-Relación.

Esta situación dificulta que los estudiantes puedan practicar con una herramienta gratuita, accesible y cercana a los ejemplos vistos en clase. En este contexto, resulta útil disponer de una aplicación que no solo permita dibujar entidades y relaciones, sino también validar el diagrama, entender su transformación al modelo relacional y observar el SQL derivado.

Por ello, la evolución de UBU E-R App tiene una motivación principalmente docente. El objetivo no es competir con herramientas profesionales de diseño de bases de datos, sino avanzar hacia una aplicación que ayude a practicar el modelado conceptual y refuerce la comprensión de las reglas de transformación al modelo relacional.

\section{Punto de partida}

El proyecto parte de una aplicación web previa, denominada Draw E-R App, orientada al modelado de diagramas Entidad-Relación con React y mxGraph. Esta versión ya ofrecía una base funcional para trabajar con entidades, atributos y relaciones, además de incluir mecanismos iniciales de validación y generación de SQL.

A partir de ese punto, este Trabajo Fin de Grado se centra en analizar el funcionamiento interno del editor, detectar limitaciones y evolucionarlo de forma incremental. El trabajo realizado no consistió en construir la herramienta desde cero, sino en mejorar la coherencia entre la representación visual del diagrama, el modelo interno, la validación, la transformación relacional, la generación SQL, la persistencia y la experiencia de edición.

Como estimación razonada, puede considerarse que aproximadamente un 40~\% del resultado final corresponde a la base heredada y un 60~\% al trabajo realizado durante este proyecto. Esta estimación no mide líneas de código, sino alcance funcional y técnico: la versión previa aportaba el editor inicial y una base de validación y generación SQL, mientras que este TFG incorpora estabilización, ampliación del soporte E/R, mejoras de transformación relacional y SQL, persistencia, usabilidad, pruebas y documentación.

La versión previa del proyecto puede consultarse en \href{https://draw-entity-relation.vercel.app/}{el despliegue original en Vercel}, mientras que la versión evolucionada está disponible en \href{https://draw-entity-relation-opal.vercel.app/}{el despliegue actual de UBU E-R App}. Las Figuras~\ref{fig:vercel_ruben} y~\ref{fig:vercel_actual} muestran la diferencia visual entre ambos despliegues: la versión heredada presenta una interfaz más básica, mientras que la versión actual incorpora una organización más completa del panel lateral, branding UBU, selector de idioma, acciones de validación, generación, importación y exportación, e información de apoyo.

\imagen{vercel_ruben}{Vista de la versión previa de Draw E-R App desplegada en Vercel.}{0.85}

\imagen{vercel_actual}{Vista de la versión evolucionada de UBU E-R App desplegada en Vercel.}{0.85}

\section{Problema abordado}

El principal problema abordado en este trabajo es la evolución de un editor de diagramas Entidad-Relación ya existente. En una herramienta de este tipo no basta con representar elementos en pantalla, ya que cada entidad, atributo o relación debe tener también un significado dentro del modelo interno de la aplicación.

Por este motivo, cualquier ampliación del editor debe tener en cuenta varios aspectos al mismo tiempo: el lienzo visual, los datos que se almacenan internamente, las reglas de validación, la transformación posterior al modelo relacional y la generación del script SQL. El trabajo se centra en avanzar en esa evolución de forma controlada, revisando el comportamiento del sistema y evitando separar la parte gráfica de la parte conceptual.

\section{Alcance del trabajo}

El alcance de este Trabajo Fin de Grado se centra en la evolución de un editor web de diagramas Entidad-Relación ya existente. El proyecto no parte de una aplicación desarrollada desde cero, sino de una herramienta heredada que ya ofrecía una base funcional para la creación de diagramas, su validación y la generación de SQL.

A partir de ese punto, el trabajo realizado incluye el análisis del sistema previo, la estabilización de partes relevantes del editor y la ampliación controlada de sus capacidades. El eje principal del desarrollo ha sido mantener la coherencia entre lo que el usuario diseña en el lienzo, la representación interna del diagrama, las validaciones aplicadas, la transformación al modelo relacional y la generación de SQL.

Dentro del alcance implementado se incluyen mejoras en la representación interna del diagrama, en las reglas de validación, en la transformación al modelo relacional y en la generación del script SQL, así como la incorporación de nuevos elementos del modelo Entidad-Relación soportados por la aplicación. También se han añadido mejoras de interacción y usabilidad orientadas a facilitar el trabajo con el editor, como la selección múltiple, el acceso a información de ayuda, opciones de importación y exportación, y ajustes visuales de la interfaz.

Al tratarse de la continuación de una aplicación existente, la contribución propia de este trabajo se delimita en el análisis del sistema heredado, la estabilización de su núcleo funcional, la ampliación de elementos E/R soportados, la revisión de la transformación relacional y SQL, las mejoras de persistencia, importación, exportación y edición, así como la documentación técnica y de usuario asociada. La base inicial del editor no se presenta como desarrollo propio completo, sino como punto de partida sobre el que se ha realizado la evolución descrita.

No se pretende construir una herramienta profesional completa de diseño de bases de datos ni cubrir todas las variantes posibles del modelo Entidad-Relación. Quedan fuera del alcance implementado funcionalidades avanzadas como las agregaciones, las restricciones completas de jerarquías ISA, la herencia múltiple, la gestión de usuarios, la persistencia remota o la configuración detallada de tipos de datos SQL. Estas limitaciones se consideran decisiones de alcance del proyecto y posibles líneas de trabajo futuro.

\section{Materiales entregados}

Junto con esta memoria se entregan los materiales necesarios para revisar, ejecutar y evaluar el proyecto:

\begin{itemize}
\item \textbf{Memoria y anexos}: documentación principal del TFG, incluyendo planificación, requisitos, diseño, documentación técnica, manual de usuario y sostenibilización curricular.

\item \textbf{Aplicación web desplegada}: versión funcional de UBU E-R App publicada en Vercel. Acceso: \href{https://draw-entity-relation-opal.vercel.app/}{aplicación desplegada en Vercel}.

\item \textbf{Repositorio Git del proyecto}: código fuente, historial de desarrollo e instrucciones de instalación y ejecución~\cite{alba2026ubuERRepo}. Acceso: \href{https://github.com/saa1002/draw-entity-relation}{repositorio del proyecto}.

\item \textbf{Material complementario}: información de ayuda disponible desde la interfaz y vídeos de demostración del funcionamiento de la aplicación.
\end{itemize}

